%!LW recipe=latexmk (xelatex)
%!TEX root = test.tex
%LTeX: enabled=false

\usepackage{unicode-math}
\setmainfont[
    BoldFont = texgyreschola-bold.otf,
    BoldItalicFont = texgyreschola-bolditalic.otf,
    ItalicFont = texgyreschola-italic.otf
]{texgyreschola-regular.otf}
\setmathfont{texgyreschola-math.otf}

\usepackage{geometry}
\geometry{
    includefoot,
    paperwidth=5.5in,
    left=0.6in,
    width=4.3in,
    paperheight=8.5in,
    top=0.6in,
    height=7.3in
    %twoside=true
}

\setlength{\parindent}{0pt}
\setlength{\parskip}{\medskipamount}

\usepackage{tabularx}

\pagestyle{plain}

% Put floats at the very top of a page.
\makeatletter
\setlength{\@fptop}{0pt}
\setlength{\@fpbot}{0pt plus 1fil}
\makeatother

\usepackage{titlesec}
\titleformat{\chapter}[display]
{\normalfont\Large\bfseries}{\chaptertitlename\ \thechapter}{10pt}{\LARGE}
\titlespacing*{\chapter}{0pt}{20pt}{20pt}

%\pdfinfoomitdate 1

%%%%%%%%%%%%%%%%%%%%%%%%%%%%%%%%%%%%%%%%%%%%%%%%%%%%%%%%%%%%%%%%%%%%%%%%%%%%%%%%

% HORIZONTAL SPACE

% The \wbox command puts its second argument in a box equal in width to its
% first argument. The alignment of the contents of the box can be specified by
% an optional argument and defaults to l.

\makeatletter

\newlength{\@wboxwidth}
\newcommand{\wbox}[3][l]{%
    \settowidth{\@wboxwidth}{#2}%
    \makebox[\@wboxwidth][#1]{#3}%
}

\makeatother

%%%%%%%%%%%%%%%%%%%%%%%%%%%%%%%%%%%%%%%%%%%%%%%%%%%%%%%%%%%%%%%%%%%%%%%%%%%%%%%%

\newcommand{\booktitle}[1]{{\itshape #1}}
\newcommand{\PSA}{\booktitle{Pocket Sky Atlas}}
\newcommand{\book}[1]{\emph{#1}}

\newcommand{\chart}[6]{
    \noindent
    \begingroup
    \raggedright
    \begin{minipage}[b][5.2cm][t]{0.50\linewidth}
        \setlength{\parskip}{\smallskipamount}
        \raggedright
        {\bfseries #1}\par
        \smallskip
        \scriptsize
        #2\par
        #3\par
        See: \ignorespaces #4\par
        #6\par
        #5\par
    \end{minipage}
    \hfill
    \input charts/#1
    \par
    \endgroup
    \chartskip

}
\newcommand{\chartskip}{\smallskip}

\newcommand{\hop}{\\→\space}
\newcommand{\degrees}[1]{\mbox{$#1^\circ$}}
\newcommand{\arcmin}[1]{\mbox{$#1'$}}
\newcommand{\hours}[1]{\mbox{$#1^{\mathrm h}$}}

\newcommand{\chartdata}[9]{
    \begin{tabularx}{\linewidth}{@{}l@{: }l@{~~}l@{: }l@{~~}l@{: }X@{}}
        RA&\wbox{\hours{00.0}}{\hours{#1}}&
        PSA&\wbox{00/00}{#4}&
        Type&#7\\
        Dec&\wbox{\degrees{+00}}{\degrees{#2}}&
        OSA&#5&
        Mag&#8\\
        Con&\wbox{Mmm}{#3}&
        CSA&\wbox{00/00/00/00}{#6}&
        Size&$#9\relax$\\
    \end{tabularx}
}

\usepackage{tikz}

\usetikzlibrary{external}
\ifinverted
\tikzexternalize[prefix=tikzexternal-inverted/]
\immediate\write18{mkdir -p tikzexternal-inverted/}
\else
\tikzexternalize[prefix=tikzexternal-correct/]
\immediate\write18{mkdir -p tikzexternal-correct/}
\fi

% We don’t define an environment for the charts since this interferes
% with the tikzexternal package.

% The chart scale in cm per degree.
\newcommand{\chartscale}{1.3}
\newcommand{\chartxscale}{
    \ifinverted
    -\chartscale
    \else
    +\chartscale
    \fi
}
\newcommand{\chartyscale}{\chartscale}
\newcommand{\starscale}{0.005}

\definecolor{chartgreen} {rgb}{0.30,0.60,0.10}
\definecolor{chartred}   {rgb}{1.00,0.00,0.00}
\definecolor{chartyellow}{rgb}{1.00,0.90,0.00}
\definecolor{chartblue}  {rgb}{0.70,0.70,1.00}

\tikzset{frame/.style={
        black,
        thick
}}
\tikzset{radialscale/.style={
        dashed,
        fill=none,
        chartblue
}}
\tikzset{compass/.style={
        black
}}

\tikzset{star/.style={
        black,
        fill=black
}}
\tikzset{OC/.style={
        black,
        fill=chartyellow,
        dotted,
        thick
}}
\tikzset{GC/.style={
        black,
        fill=chartyellow,
        thick
}}
\tikzset{PN/.style={
        black,
        fill=chartgreen,
        thick
}}
\tikzset{BN/.style={
        black,
        fill=chartgreen,
        thick
}}
\tikzset{DN/.style={
        black,
        fill=none,
        thick,
        dashed
}}
\tikzset{SC/.style={
        black,
        fill=none,
        thick,
        dotted
}}
\tikzset{AST/.style={
        black,
        fill=none,
        thick,
        dotted
}}
\tikzset{DS/.style={
        black,
        fill=none,
        thick
}}
\tikzset{GAL/.style={
        black,
        fill=chartred,
        thick
}}

\newcommand{\drawframe}[2]{
    \draw [frame] (-#2,-#2) -- (-#2,+#2) -- (+#2,+#2) -- (+#2,-#2) -- cycle;
    \clip (-#2,-#2) -- (-#2,+#2) -- (+#2,+#2) -- (+#2,-#2) -- cycle;
}
\newcommand{\drawradialscale}[1]{
    \filldraw [radialscale] (0,0) circle [radius=#1];
}
\newcommand{\drawcompass}[3]{
    \draw[compass,->] (#1,#2) -- +(0.0,#3) node [anchor=south] {\tiny N};
    \ifinverted
    \draw[compass,->] (#1,#2) -- +(-#3,0.0) node [anchor=west] {\tiny E};
    \else
    \draw[compass,->] (#1,#2) -- +(-#3,0.0) node [anchor=east] {\tiny E};
    \fi
}
\newcommand{\drawstar}[3]{
    \filldraw [star] (#1,#2) circle [radius=(#3-10.5)*\starscale];
}
\newcommand{\drawOC}[3]{
    \filldraw [OC] (#1,#2) circle [radius=#3];
}
\newcommand{\drawGC}[3]{
    \filldraw [GC] (#1,#2) circle [radius=#3];
    \draw [GC] (#1,#2) --   +(0:#3);
    \draw [GC] (#1,#2) --  +(90:#3);
    \draw [GC] (#1,#2) -- +(180:#3);
    \draw [GC] (#1,#2) -- +(270:#3);
}
\newcommand{\drawPN}[3]{
    \filldraw [PN] (#1,#2) circle [radius=#3];
    \draw [PN] (#1,#2) ++   (0:#3) --   +(0:#3);
    \draw [PN] (#1,#2) ++  (90:#3) --  +(90:#3);
    \draw [PN] (#1,#2) ++ (180:#3) -- +(180:#3);
    \draw [PN] (#1,#2) ++ (270:#3) -- +(270:#3);
}
\newcommand{\drawBN}[3]{
    \filldraw [BN] (#1,#2) circle [radius=#3];
}
\newcommand{\drawDN}[3]{
    \filldraw [DN] (#1,#2) circle [radius=#3];
}
\newcommand{\drawSC}[5]{
    \filldraw [SC] (#1,#2) circle [x radius = #4, y radius=#3, rotate=#5];
}
\newcommand{\drawAST}[5]{
    \filldraw [AST] (#1,#2) circle [x radius = #4, y radius=#3, rotate=#5];
}
\newcommand{\drawGAL}[5]{
    \filldraw [GAL] (#1,#2) circle [x radius = #4, y radius=#3, rotate=#5];
}
\newcommand{\drawDS}[3]{
    \draw [DS] (#1,#2) --  +(45:#3);
    \draw [DS] (#1,#2) -- +(135:#3);
    \draw [DS] (#1,#2) -- +(225:#3);
    \draw [DS] (#1,#2) -- +(315:#3);
}
